%
\begin{isabellebody}%
\setisabellecontext{Week{\isacharunderscore}{\kern0pt}{\isadigit{2}}{\isacharunderscore}{\kern0pt}sol}%
%
\isadelimtheory
\isanewline
%
\endisadelimtheory
%
\isatagtheory
\isacommand{theory}\isamarkupfalse%
\ Week{\isacharunderscore}{\kern0pt}{\isadigit{2}}{\isacharunderscore}{\kern0pt}sol\isanewline
\isakeyword{imports}\ Main\isanewline
\isakeyword{begin}%
\endisatagtheory
{\isafoldtheory}%
%
\isadelimtheory
%
\endisadelimtheory
%
\begin{isamarkuptext}%
\vspace{15ex}\ExerciseSheet{7}{}%
\end{isamarkuptext}\isamarkuptrue%
%
\begin{isamarkuptext}%
\Exercise{Fold function}
  The fold function is a very generic function, that can be used to express 
  multiple other interesting functions over lists.%
\end{isamarkuptext}\isamarkuptrue%
%
\begin{isamarkuptext}%
Have a look at Isabelle/HOL's standard function \isa{fold}, defined as follows:%
\end{isamarkuptext}\isamarkuptrue%
\isacommand{fun}\isamarkupfalse%
\ fold{\isacharcolon}{\kern0pt}{\isacharcolon}{\kern0pt}{\isachardoublequoteopen}{\isacharparenleft}{\kern0pt}{\isacharprime}{\kern0pt}a\ {\isasymRightarrow}\ {\isacharprime}{\kern0pt}b\ {\isasymRightarrow}\ {\isacharprime}{\kern0pt}b{\isacharparenright}{\kern0pt}\ {\isasymRightarrow}\ {\isacharprime}{\kern0pt}a\ list\ {\isasymRightarrow}\ {\isacharprime}{\kern0pt}b\ {\isasymRightarrow}\ {\isacharprime}{\kern0pt}b{\isachardoublequoteclose}\ \isakeyword{where}\isanewline
{\isachardoublequoteopen}fold\ f\ {\isacharbrackleft}{\kern0pt}{\isacharbrackright}{\kern0pt}\ acc\ {\isacharequal}{\kern0pt}\ acc{\isachardoublequoteclose}\isanewline
{\isacharbar}{\kern0pt}\ {\isachardoublequoteopen}fold\ f\ {\isacharparenleft}{\kern0pt}x\ {\isacharhash}{\kern0pt}xs{\isacharparenright}{\kern0pt}\ acc\ {\isacharequal}{\kern0pt}\ fold\ f\ xs\ {\isacharparenleft}{\kern0pt}f\ x\ acc{\isacharparenright}{\kern0pt}{\isachardoublequoteclose}%
\begin{isamarkuptext}%
Write a function to compute the sum of the elements of a list. 
  Define two versions, one direct recursive definition, and one using fold.
  Show that both are equal.

  Hint: use automation!%
\end{isamarkuptext}\isamarkuptrue%
\isacommand{fun}\isamarkupfalse%
\ list{\isacharunderscore}{\kern0pt}sum\ {\isacharcolon}{\kern0pt}{\isacharcolon}{\kern0pt}\ {\isachardoublequoteopen}nat\ list\ {\isasymRightarrow}\ nat{\isachardoublequoteclose}\ \isanewline
\ \ \ \ \isanewline
\ \ \isakeyword{where}\isanewline
\ \ \ \ {\isachardoublequoteopen}list{\isacharunderscore}{\kern0pt}sum\ {\isacharbrackleft}{\kern0pt}{\isacharbrackright}{\kern0pt}\ {\isacharequal}{\kern0pt}\ {\isadigit{0}}{\isachardoublequoteclose}\isanewline
\ \ {\isacharbar}{\kern0pt}\ {\isachardoublequoteopen}list{\isacharunderscore}{\kern0pt}sum\ {\isacharparenleft}{\kern0pt}x{\isacharhash}{\kern0pt}xs{\isacharparenright}{\kern0pt}\ {\isacharequal}{\kern0pt}\ x\ {\isacharplus}{\kern0pt}\ list{\isacharunderscore}{\kern0pt}sum\ xs{\isachardoublequoteclose}\ \ \isanewline
\ \ \ \ \ \ \isanewline
\isanewline
\isacommand{definition}\isamarkupfalse%
\ list{\isacharunderscore}{\kern0pt}sum{\isacharprime}{\kern0pt}\ {\isacharcolon}{\kern0pt}{\isacharcolon}{\kern0pt}\ {\isachardoublequoteopen}nat\ list\ {\isasymRightarrow}\ nat{\isachardoublequoteclose}\isanewline
\ \ \ \ \isanewline
\ \ \isakeyword{where}\ {\isachardoublequoteopen}list{\isacharunderscore}{\kern0pt}sum{\isacharprime}{\kern0pt}\ xs\ {\isasymequiv}\ fold\ {\isacharparenleft}{\kern0pt}{\isacharplus}{\kern0pt}{\isacharparenright}{\kern0pt}\ xs\ {\isadigit{0}}{\isachardoublequoteclose}\isanewline
\ \ \ \ \ \ \ \ \isanewline
\isanewline
\ \ \ \ \isanewline
\isacommand{lemma}\isamarkupfalse%
\ aux{\isacharcolon}{\kern0pt}\ {\isachardoublequoteopen}fold\ {\isacharparenleft}{\kern0pt}{\isacharplus}{\kern0pt}{\isacharparenright}{\kern0pt}\ xs\ a\ {\isacharequal}{\kern0pt}\ list{\isacharunderscore}{\kern0pt}sum\ xs\ {\isacharplus}{\kern0pt}\ a{\isachardoublequoteclose}\ \ \isanewline
%
\isadelimproof
\ \ %
\endisadelimproof
%
\isatagproof
\isacommand{by}\isamarkupfalse%
\ {\isacharparenleft}{\kern0pt}induction\ xs\ arbitrary{\isacharcolon}{\kern0pt}\ a{\isacharparenright}{\kern0pt}\ auto%
\endisatagproof
{\isafoldproof}%
%
\isadelimproof
\isanewline
%
\endisadelimproof
\ \ \ \ \ \ \ \ \isanewline
\isanewline
\isacommand{lemma}\isamarkupfalse%
\ {\isachardoublequoteopen}list{\isacharunderscore}{\kern0pt}sum\ xs\ {\isacharequal}{\kern0pt}\ list{\isacharunderscore}{\kern0pt}sum{\isacharprime}{\kern0pt}\ xs{\isachardoublequoteclose}\isanewline
%
\isadelimproof
\ \ \ \ \ \ \isanewline
\ \ %
\endisadelimproof
%
\isatagproof
\isacommand{unfolding}\isamarkupfalse%
\ list{\isacharunderscore}{\kern0pt}sum{\isacharprime}{\kern0pt}{\isacharunderscore}{\kern0pt}def\ \isacommand{by}\isamarkupfalse%
\ {\isacharparenleft}{\kern0pt}simp\ add{\isacharcolon}{\kern0pt}\ aux{\isacharparenright}{\kern0pt}%
\endisatagproof
{\isafoldproof}%
%
\isadelimproof
%
\endisadelimproof
%
\begin{isamarkuptext}%
\Exercise{Tail recursive \isa{reverse}}
In Isabelle, there is the reverse function \isa{rev}, which, given a list, computes another list
with the same elements but in reverse order.
Take a loot into its definition:%
\end{isamarkuptext}\isamarkuptrue%
\isacommand{thm}\isamarkupfalse%
\ rev{\isachardot}{\kern0pt}simps%
\begin{isamarkuptext}%
Recall tail recursion: a function is tail recursive if in all recursive calls it does not call 
  anything after itself\footnote{See \url{https://en.wikipedia.org/wiki/Tail_call.} for a quick review.}.
  For instance, the implementaion of \isa{rev} in Isabelle is not tail recursive, because it calls
  \isa{{\isacharparenleft}{\kern0pt}{\isacharat}{\kern0pt}{\isacharparenright}{\kern0pt}} after it calls \isa{rev} in the second equation.
  
  Write an alternative implementation of the reverse function that is tail recursive.

  Hint: recall that you can most of the times derive a tail recursive function using an accumulator
        argument.%
\end{isamarkuptext}\isamarkuptrue%
\isacommand{fun}\isamarkupfalse%
\ rev{\isacharprime}{\kern0pt}{\isacharcolon}{\kern0pt}{\isacharcolon}{\kern0pt}{\isachardoublequoteopen}{\isacharprime}{\kern0pt}a\ list\ {\isasymRightarrow}\ {\isacharprime}{\kern0pt}a\ list\ {\isasymRightarrow}\ {\isacharprime}{\kern0pt}a\ list{\isachardoublequoteclose}\ \isakeyword{where}\isanewline
\ \ {\isachardoublequoteopen}rev{\isacharprime}{\kern0pt}\ {\isacharbrackleft}{\kern0pt}{\isacharbrackright}{\kern0pt}\ acc\ {\isacharequal}{\kern0pt}\ acc{\isachardoublequoteclose}\isanewline
{\isacharbar}{\kern0pt}\ {\isachardoublequoteopen}rev{\isacharprime}{\kern0pt}\ {\isacharparenleft}{\kern0pt}x\ {\isacharhash}{\kern0pt}xs{\isacharparenright}{\kern0pt}\ acc\ {\isacharequal}{\kern0pt}\ rev{\isacharprime}{\kern0pt}\ xs\ {\isacharparenleft}{\kern0pt}x\ {\isacharhash}{\kern0pt}\ acc{\isacharparenright}{\kern0pt}{\isachardoublequoteclose}\isanewline
\isanewline
\isanewline
\isanewline
\isanewline
\isacommand{fun}\isamarkupfalse%
\ rev{\isacharunderscore}{\kern0pt}tr{\isacharcolon}{\kern0pt}{\isacharcolon}{\kern0pt}{\isachardoublequoteopen}{\isacharprime}{\kern0pt}a\ list\ {\isasymRightarrow}\ {\isacharprime}{\kern0pt}a\ list{\isachardoublequoteclose}\ \isakeyword{where}\isanewline
\isanewline
\ \ {\isachardoublequoteopen}rev{\isacharunderscore}{\kern0pt}tr\ xs\ {\isacharequal}{\kern0pt}\ rev{\isacharprime}{\kern0pt}\ xs\ {\isacharbrackleft}{\kern0pt}{\isacharbrackright}{\kern0pt}{\isachardoublequoteclose}%
\begin{isamarkuptext}%
Show that the two implementations are equivalent.%
\end{isamarkuptext}\isamarkuptrue%
\isacommand{lemma}\isamarkupfalse%
\ rev{\isacharprime}{\kern0pt}{\isacharunderscore}{\kern0pt}append{\isacharcolon}{\kern0pt}\ {\isachardoublequoteopen}rev{\isacharprime}{\kern0pt}\ xs\ ys\ {\isacharequal}{\kern0pt}\ {\isacharparenleft}{\kern0pt}rev{\isacharprime}{\kern0pt}\ xs\ {\isacharbrackleft}{\kern0pt}{\isacharbrackright}{\kern0pt}{\isacharparenright}{\kern0pt}\ {\isacharat}{\kern0pt}\ ys{\isachardoublequoteclose}\isanewline
%
\isadelimproof
%
\endisadelimproof
%
\isatagproof
\isacommand{proof}\isamarkupfalse%
\ {\isacharparenleft}{\kern0pt}induction\ xs\ arbitrary{\isacharcolon}{\kern0pt}\ ys{\isacharparenright}{\kern0pt}\isanewline
\ \ \isacommand{case}\isamarkupfalse%
\ Nil\isanewline
\ \ \isacommand{then}\isamarkupfalse%
\ \isacommand{show}\isamarkupfalse%
\ {\isacharquery}{\kern0pt}case\ \isanewline
\ \ \ \ \isacommand{by}\isamarkupfalse%
\ auto\isanewline
\isacommand{next}\isamarkupfalse%
\isanewline
\ \ \isacommand{case}\isamarkupfalse%
\ {\isacharparenleft}{\kern0pt}Cons\ a\ xs{\isacharparenright}{\kern0pt}\isanewline
\ \ \isacommand{have}\isamarkupfalse%
\ {\isadigit{1}}{\isacharcolon}{\kern0pt}\ {\isachardoublequoteopen}rev{\isacharprime}{\kern0pt}\ xs\ {\isacharparenleft}{\kern0pt}a\ {\isacharhash}{\kern0pt}\ ys{\isacharparenright}{\kern0pt}\ {\isacharequal}{\kern0pt}\ {\isacharparenleft}{\kern0pt}rev{\isacharprime}{\kern0pt}\ xs\ {\isacharbrackleft}{\kern0pt}{\isacharbrackright}{\kern0pt}{\isacharparenright}{\kern0pt}\ {\isacharat}{\kern0pt}\ {\isacharparenleft}{\kern0pt}a\ {\isacharhash}{\kern0pt}\ ys{\isacharparenright}{\kern0pt}{\isachardoublequoteclose}\isanewline
\ \ \ \ \isacommand{apply}\isamarkupfalse%
{\isacharparenleft}{\kern0pt}subst\ Cons{\isacharparenright}{\kern0pt}\isanewline
\ \ \ \ \isacommand{{\isachardot}{\kern0pt}{\isachardot}{\kern0pt}}\isamarkupfalse%
\isanewline
\ \ \isacommand{have}\isamarkupfalse%
\ {\isadigit{2}}{\isacharcolon}{\kern0pt}\ {\isachardoublequoteopen}rev{\isacharprime}{\kern0pt}\ xs\ {\isacharbrackleft}{\kern0pt}a{\isacharbrackright}{\kern0pt}\ {\isacharequal}{\kern0pt}\ {\isacharparenleft}{\kern0pt}rev{\isacharprime}{\kern0pt}\ xs\ {\isacharbrackleft}{\kern0pt}{\isacharbrackright}{\kern0pt}{\isacharparenright}{\kern0pt}\ {\isacharat}{\kern0pt}\ {\isacharbrackleft}{\kern0pt}a{\isacharbrackright}{\kern0pt}{\isachardoublequoteclose}\isanewline
\ \ \ \ \isacommand{apply}\isamarkupfalse%
{\isacharparenleft}{\kern0pt}subst\ Cons{\isacharparenright}{\kern0pt}\isanewline
\ \ \ \ \isacommand{{\isachardot}{\kern0pt}{\isachardot}{\kern0pt}}\isamarkupfalse%
\isanewline
\ \ \isacommand{show}\isamarkupfalse%
\ {\isacharquery}{\kern0pt}case\isanewline
\ \ \ \ \isacommand{using}\isamarkupfalse%
\ {\isadigit{1}}\ {\isadigit{2}}\isanewline
\ \ \ \isacommand{by}\isamarkupfalse%
\ simp\isanewline
\isacommand{qed}\isamarkupfalse%
%
\endisatagproof
{\isafoldproof}%
%
\isadelimproof
\isanewline
%
\endisadelimproof
\isanewline
\ \isanewline
\isacommand{lemma}\isamarkupfalse%
\ {\isachardoublequoteopen}rev{\isacharunderscore}{\kern0pt}tr\ xs\ {\isacharequal}{\kern0pt}\ rev\ xs{\isachardoublequoteclose}\isanewline
%
\isadelimproof
\isanewline
%
\endisadelimproof
%
\isatagproof
\isacommand{proof}\isamarkupfalse%
{\isacharparenleft}{\kern0pt}induction\ xs{\isacharparenright}{\kern0pt}\isanewline
\ \ \isacommand{case}\isamarkupfalse%
\ Nil\isanewline
\ \ \isacommand{then}\isamarkupfalse%
\ \isacommand{show}\isamarkupfalse%
\ {\isacharquery}{\kern0pt}case\isanewline
\ \ \ \ \isacommand{by}\isamarkupfalse%
\ auto\isanewline
\isacommand{next}\isamarkupfalse%
\isanewline
\ \ \isacommand{case}\isamarkupfalse%
\ {\isacharparenleft}{\kern0pt}Cons\ a\ xs{\isacharparenright}{\kern0pt}\isanewline
\ \ \isacommand{then}\isamarkupfalse%
\ \isacommand{show}\isamarkupfalse%
\ {\isacharquery}{\kern0pt}case\isanewline
\ \ \ \ \isacommand{apply}\isamarkupfalse%
\ simp\isanewline
\ \ \ \ \isacommand{apply}\isamarkupfalse%
{\isacharparenleft}{\kern0pt}subst\ rev{\isacharprime}{\kern0pt}{\isacharunderscore}{\kern0pt}append{\isacharparenright}{\kern0pt}\isanewline
\ \ \ \ \isacommand{by}\isamarkupfalse%
\ auto\isanewline
\isacommand{qed}\isamarkupfalse%
%
\endisatagproof
{\isafoldproof}%
%
\isadelimproof
\isanewline
%
\endisadelimproof
\isanewline
\isanewline
%
\isadelimtheory
\isanewline
%
\endisadelimtheory
%
\isatagtheory
\isacommand{end}\isamarkupfalse%
%
\endisatagtheory
{\isafoldtheory}%
%
\isadelimtheory
%
\endisadelimtheory
%
\end{isabellebody}%
\endinput
%:%file=Week_2_sol.tex%:%
%:%6=1%:%
%:%11=2%:%
%:%12=2%:%
%:%13=3%:%
%:%14=4%:%
%:%23=6%:%
%:%27=8%:%
%:%28=9%:%
%:%29=10%:%
%:%33=13%:%
%:%35=15%:%
%:%36=15%:%
%:%37=16%:%
%:%38=17%:%
%:%40=20%:%
%:%41=21%:%
%:%42=22%:%
%:%43=23%:%
%:%44=24%:%
%:%46=29%:%
%:%47=29%:%
%:%48=30%:%
%:%49=31%:%
%:%50=32%:%
%:%51=33%:%
%:%52=34%:%
%:%53=35%:%
%:%54=36%:%
%:%55=36%:%
%:%56=37%:%
%:%57=38%:%
%:%58=39%:%
%:%59=40%:%
%:%60=41%:%
%:%61=42%:%
%:%62=42%:%
%:%65=43%:%
%:%69=43%:%
%:%70=43%:%
%:%75=43%:%
%:%78=44%:%
%:%79=45%:%
%:%80=46%:%
%:%81=46%:%
%:%84=47%:%
%:%85=48%:%
%:%89=48%:%
%:%90=48%:%
%:%91=48%:%
%:%100=51%:%
%:%101=52%:%
%:%102=53%:%
%:%103=54%:%
%:%105=57%:%
%:%106=57%:%
%:%108=60%:%
%:%109=61%:%
%:%110=62%:%
%:%111=63%:%
%:%112=64%:%
%:%113=65%:%
%:%114=66%:%
%:%115=67%:%
%:%116=68%:%
%:%118=73%:%
%:%119=73%:%
%:%120=74%:%
%:%121=75%:%
%:%122=76%:%
%:%123=77%:%
%:%124=78%:%
%:%125=79%:%
%:%126=80%:%
%:%127=80%:%
%:%128=81%:%
%:%129=82%:%
%:%131=85%:%
%:%133=90%:%
%:%134=90%:%
%:%141=91%:%
%:%142=91%:%
%:%143=92%:%
%:%144=92%:%
%:%145=93%:%
%:%146=93%:%
%:%147=93%:%
%:%148=94%:%
%:%149=94%:%
%:%150=95%:%
%:%151=95%:%
%:%152=96%:%
%:%153=96%:%
%:%154=97%:%
%:%155=97%:%
%:%156=98%:%
%:%157=98%:%
%:%158=99%:%
%:%159=99%:%
%:%160=100%:%
%:%161=100%:%
%:%162=101%:%
%:%163=101%:%
%:%164=102%:%
%:%165=102%:%
%:%166=103%:%
%:%167=103%:%
%:%168=104%:%
%:%169=104%:%
%:%170=105%:%
%:%171=105%:%
%:%172=106%:%
%:%178=106%:%
%:%181=107%:%
%:%182=108%:%
%:%183=109%:%
%:%184=109%:%
%:%187=110%:%
%:%192=111%:%
%:%193=111%:%
%:%194=112%:%
%:%195=112%:%
%:%196=113%:%
%:%197=113%:%
%:%198=113%:%
%:%199=114%:%
%:%200=114%:%
%:%201=115%:%
%:%202=115%:%
%:%203=116%:%
%:%204=116%:%
%:%205=117%:%
%:%206=117%:%
%:%207=117%:%
%:%208=118%:%
%:%209=118%:%
%:%210=119%:%
%:%211=119%:%
%:%212=120%:%
%:%213=120%:%
%:%214=121%:%
%:%220=121%:%
%:%223=122%:%
%:%224=123%:%
%:%227=124%:%
%:%232=125%:%
